\documentclass[letterpaper,11pt]{article}

\usepackage{amsmath, amsfonts, amssymb, amsthm}
\usepackage{lmodern}
\usepackage{enumitem}
\usepackage{hyperref}
\usepackage[x11names, rgb]{xcolor}
\usepackage{graphicx}
\usepackage{fullpage}
\usepackage{etoolbox}
\usepackage{mathtools}


\newcommand{\cX}{\mathcal{X}}
\newcommand{\cC}{\mathcal{C}}
\newcommand{\cF}{\mathcal{F}}
\newcommand{\R}{\mathbb{R}}
\newcommand{\dd}{{\rm d}}
\DeclareMathOperator{\sech}{sech}


\setlength{\parindent}{0pt}
\setlength{\parskip}{6pt}

\newcommand{\nopagenumbers}{
    \pagestyle{empty}
}

\newenvironment{indented}{
    \begin{list}{}{\leftmargin=1em}
    \item[]
}{
    \end{list}
}

\newtoggle{questionnumbers}
\toggletrue{questionnumbers}

\newcounter{questionCounter}
\newcounter{partCounter}[questionCounter]

\newcommand{\question}[1]{%
    \stepcounter{questionCounter}%
    \setcounter{partCounter}{0}%
    \vspace{0.25in}
    \hrule
    \vspace{0.4em}
    \noindent\textbf{Question \thequestionCounter:\ }#1
    \vspace{0.8em}
    \hrule
    \vspace{0.1in}
}

\newcommand{\qpart}{%
    \stepcounter{partCounter}%
    \begin{indented}
    \textbf{(\alph{partCounter}) }
}

\newcommand{\closepart}{
    \end{indented}
}

\newcommand{\myhwname}{AMATH 353 HW 7}
\newcommand{\myname}{Jordan Tucker}

\newcommand{\header}{
    \begin{center}
    {\Large \myhwname}\\
    \myname\\
    \today
    \end{center}
}

\nopagenumbers
% Document
\begin{document}
    \header
    \question{
        \textbf{An Easier Shock Solution}\\
        One setting where the method of characteristics fails is when
        the characteristics intersect and the solution becomes infinitely steep.
        This leads to the formation of \textit{shock wave}, and
        we saw that the shock moves along a characteristic
        $(x_s(t),t)$ which can be determined from Rankine-Hugoniot condition
        \[
            \frac{dx_s(t)}{dt} =\frac{\phi(u^+)-\phi(u^-)}{u^+-u^-}.
        \]
        Enforcing the Rankine-Hugoniot condition, determine
        the shock solution of the IVP
        \[
            \begin{cases}
                u_t + uu_x = 0, & x \in \R, t \ge 0,\\
                u(x,0)= u_0(x), & x \in \R,
            \end{cases}
            \quad \text{where}\quad
            u_0(x) =
            \begin{cases}
                2, &x<-1,\\
                1-x, & |x|\le 1.\\
                0, &x>1.\\
            \end{cases}
        \]
    }
    Find the characteristic curve.
    \begin{gather*}
        u_t + uu_x =0\\
        c(u) = u\\
        \frac{dx}{dt} = u\\
        u = u_0(x_0)\\
        \frac{dx}{dt} = u\\
        dx = u dt\\
        \int dx = \int u \,dt\\
        x(t) = u_0(x_0) t + C\\
    \end{gather*}
    Solving for $C$
    \begin{gather*}
        x(0) = x_0\\
        x_0 = u_0(x_0) (0) + C\\
        x_0 = C\\\\
        x(t) = u_0(x_0) t + x_0
    \end{gather*}
    For the first region, $x_0<-1,\ u_0(x_0)=2$
    \[
        x(t) = 2t +x_0
    \]
    For the second region $x_0 > 0, \ u_0 =0$
    \[
        x(t) = (0) t +x_0 = X_0
    \]
    For the third region where the break occurs $-1 \leq x \leq 1,\ u_0(x_0)=1-x_0$
    \begin{gather*}
        u = 1-x_0\\
        x(t) = (1-x_0)t + x_0\\
        x_0 = \frac{x-t}{1-t}\\
        u(x,t) = \frac{1-x}{1-t}\\
        t_s = 1
    \end{gather*}
    \begin{gather*}
        \phi(u) = \frac{u^2}{2}\\
        u^- = 2, \ u^+ = 0\\
        \frac{ds}{dt} = \frac{0-2}{0-2} = 1\\
        x_s(t) = t
    \end{gather*}
    Final piece wise solution:
    \begin{gather*}
        \text{For $t \leq 1$}\\
        u(x,t) = \begin{cases}
                     2,\ x<2t-1\\
                     \frac{1-x}{1-t},\ 2t-1\leq x \leq 1\\
                     0,\ x>1
        \end{cases}\\\\
        \text{For $t>1$}\\
        u(x,t) = \begin{cases}
            2,\ x< t\\
                     0,\ x>t
            \end{cases}
    \end{gather*}
    \newpage
    \question{
        \textbf{A Harder Shock Solution}\\
        Enforcing the Rankine-Hugoniot condition, determine
        the shock solution of the IVP
        \[
            \begin{cases}
                u_t + uu_x = 0, & x \in \R, t \ge 0,\\
                u(x,0)= u_0(x), & x \in \R,
            \end{cases}
            \quad \text{where}\quad
            u_0(x) =
            \begin{cases}
                0, &x<0,\\
                x, & 0\le x\le 1.\\
                0, &1<x.\\
            \end{cases}
        \]

    }

    For the first region, $x_0 < 0,\ u_0(x_0) = 0$
    \begin{gather*}
        x(t) = (0)t + x_0 \\
        x = x_0 \\
        u(x,t) = 0
    \end{gather*}
    For the second region, $0 \leq x_0 \leq 1,\ u_0(x_0) = x_0$
    \begin{gather*}
        x(t) = x_0 t + x_0 = x_0(t + 1) \\
        x_0 = \frac{x}{t + 1} \\
        u(x,t) = \frac{x}{t + 1}
    \end{gather*}
    For the third region, $x_0 > 1,\ u_0(x_0) = 0$
    \begin{gather*}
        x(t) = (0)t + x_0 \\
        x = x_0 \\
        u(x,t) = 0
    \end{gather*}
    Taking the derivative of the $x_s$ with respect to $t$ to find the shock speed between region 2 and 3:
    \begin{gather*}
        t_s = 0,\ x_s(0) = 1 \\
        \frac{dx_s}{dt} = \frac{u^+ + u^-}{2} = \frac{0 + \frac{x_s}{1+t}}{2} = \frac{x_s}{2+2t}
    \end{gather*}
    Solving the separable differential equation:
    \begin{gather*}
        \frac{dx_s}{x_s} = \frac{dt}{2(1+t)} \\
        \int \frac{1}{x_s} \,dx_s = \int \frac{1}{2(1+t)} \,dt \\
        \ln|x_s| = \frac{1}{2}\ln|1+t| + C
    \end{gather*}
    Using the initial condition $x_s(0) = 1$:
    \begin{gather*}
        \ln(1) = \frac{1}{2}\ln(1) + C \\
        0 = 0 + C \\
        C = 0
    \end{gather*}
    \begin{gather*}
        \ln|x_s| = \ln\left(\sqrt{1+t}\right) \\
        x_s(t) = \sqrt{t+1}
    \end{gather*}
    Final piecewise solution:
    \begin{gather*}
        u(x,t) = \begin{cases}
                     0, & x < 0 \\
                     \frac{x}{t+1}, & 0 \leq x < \sqrt{t+1} \\
                     0, & x > \sqrt{t+1}
        \end{cases}
    \end{gather*}
    \newpage
    \question{
        \textbf{Three Shocks}\\
        Enforcing the Rankine-Hugoniot condition, determine
        the shock solution of the IVP
        \[
            \begin{cases}
                u_t + u^2u_x = 0, & x \in \R, t \ge 0,\\
                u(x,0)= u_0(x), & x \in \R,
            \end{cases}
            \quad \text{where}\quad
            u_0(x) =
            \begin{cases}
                2, &x<0,\\
                1, & 0\le x\le 1.\\
                0, &1<x.\\
            \end{cases}
        \]
    }

    Finding the flux function $\phi(u)$ by taking the integral:
    \begin{gather*}
        u_t + \phi'(u) u_x = 0 \\
        \phi'(u) = u^2 \\
        \phi(u) = \int u^2 \,du = \frac{1}{3}u^3
    \end{gather*}
    For the first shock at $x_0 = 0$, taking the derivative of $x_{s1}$:
    \begin{gather*}
        u^- = 2,\ u^+ = 1 \\
        \frac{dx_{s1}}{dt} = \frac{\phi(u^+) - \phi(u^-)}{u^+ - u^-} = \frac{\frac{1}{3}(1)^3 - \frac{1}{3}(2)^3}{1 - 2} = \frac{\frac{1}{3} - \frac{8}{3}}{-1} = \frac{7}{3} \\
        dx_{s1} = \frac{7}{3} \,dt \\
        \int \,dx_{s1} = \int \frac{7}{3} \,dt \\
        x_{s1}(t) = \frac{7}{3}t + C
    \end{gather*}
    Using $x_{s1}(0) = 0$:
    \begin{gather*}
        0 = \frac{7}{3}(0) + C \implies C = 0 \\
        x_{s1}(t) = \frac{7}{3}t
    \end{gather*}
    For the second shock at $x_0 = 1$, taking the derivative of $x_{s2}$:
    \begin{gather*}
        u^- = 1,\ u^+ = 0 \\
        \frac{dx_{s2}}{dt} = \frac{\frac{1}{3}(0)^3 - \frac{1}{3}(1)^3}{0 - 1} = \frac{-\frac{1}{3}}{-1} = \frac{1}{3} \\
        dx_{s2} = \frac{1}{3} \,dt \\
        \int \,dx_{s2} = \int \frac{1}{3} \,dt \\
        x_{s2}(t) = \frac{1}{3}t + C
    \end{gather*}
    Using $x_{s2}(0) = 1$:
    \begin{gather*}
        1 = \frac{1}{3}(0) + C \implies C = 1 \\
        x_{s2}(t) = \frac{1}{3}t + 1
    \end{gather*}
    Solving for the collision time $t_c$:
    \begin{gather*}
        \frac{7}{3}t = \frac{1}{3}t + 1 \\
        7t = t + 3 \\
        6t = 3 \\
        t_c = \frac{1}{2} \\
        x_c = x_{s1}\left(\frac{1}{2}\right) = \frac{7}{6}
    \end{gather*}
    For the third shock, taking the derivative of $x_{s3}$:
    \begin{gather*}
        u^- = 2,\ u^+ = 0 \\
        \frac{dx_{s3}}{dt} = \frac{\frac{1}{3}(0)^3 - \frac{1}{3}(2)^3}{0 - 2} = \frac{-\frac{8}{3}}{-2} = \frac{4}{3} \\
        dx_{s3} = \frac{4}{3} \,dt \\
        \int \,dx_{s3} = \int \frac{4}{3} \,dt \\
        x_{s3}(t) = \frac{4}{3}t + C
    \end{gather*}
    Using the point $(x_c, t_c) = \left(\frac{7}{6}, \frac{1}{2}\right)$:
    \begin{gather*}
        \frac{7}{6} = \frac{4}{3}\left(\frac{1}{2}\right) + C \\
        \frac{7}{6} = \frac{4}{6} + C \implies C = \frac{3}{6} = \frac{1}{2} \\
        x_{s3}(t) = \frac{4}{3}t + \frac{1}{2}
    \end{gather*}
    Final piecewise solution:
    \begin{gather*}
        \text{For } t \leq \frac{1}{2} \\
        u(x,t) = \begin{cases}
                     2, & x < \frac{7}{3}t \\
                     1, & \frac{7}{3}t \leq x \leq \frac{1}{3}t + 1 \\
                     0, & x > \frac{1}{3}t + 1
        \end{cases} \\\\
        \text{For } t > \frac{1}{2} \\
        u(x,t) = \begin{cases}
                     2, & x < \frac{4}{3}t + \frac{1}{2} \\
                     0, & x > \frac{4}{3}t + \frac{1}{2}
        \end{cases}
    \end{gather*}


    \newpage
    \question{
        \textbf{Two shocks}
        Enforcing the Rankine-Hugoniot condition, determine
        the shock solution of the IVP
        \[
            \begin{cases}
                u_t + uu_x = 0, & x \in \R, t \ge 0,\\
                u(x,0)= u_0(x), & x \in \R,
            \end{cases}
            \quad \text{where}\quad
            u_0(x) =
            \begin{cases}
                1, &x<-1,\\
                0, & -1\le x\le 1.\\
                -1, &1<x.\\
            \end{cases}
        \]
    }

    \begin{gather*}
        u_t + u u_x = 0 \\
        \phi(u) = \frac{1}{2}u^2
    \end{gather*}
    For the first shock at $x_0 = -1$, taking the derivative of $x_{s1}$:
    \begin{gather*}
        u^- = 1,\ u^+ = 0 \\
        \frac{dx_{s1}}{dt} = \frac{\frac{1}{2}(0)^2 - \frac{1}{2}(1)^2}{0 - 1} = \frac{-\frac{1}{2}}{-1} = \frac{1}{2} \\
        dx_{s1} = \frac{1}{2} \,dt \\
        \int \,dx_{s1} = \int \frac{1}{2} \,dt \\
        x_{s1}(t) = \frac{1}{2}t + C
    \end{gather*}
    Using $x_{s1}(0) = -1$:
    \begin{gather*}
        -1 = \frac{1}{2}(0) + C \implies C = -1 \\
        x_{s1}(t) = \frac{1}{2}t - 1
    \end{gather*}
    For the second shock at $x_0 = 1$, taking the derivative of $x_{s2}$:
    \begin{gather*}
        u^- = 0,\ u^+ = -1 \\
        \frac{dx_{s2}}{dt} = \frac{\frac{1}{2}(-1)^2 - \frac{1}{2}(0)^2}{-1 - 0} = \frac{\frac{1}{2}}{-1} = -\frac{1}{2} \\
        dx_{s2} = -\frac{1}{2} \,dt \\
        \int \,dx_{s2} = \int -\frac{1}{2} \,dt \\
        x_{s2}(t) = -\frac{1}{2}t + C
    \end{gather*}
    Using $x_{s2}(0) = 1$:
    \begin{gather*}
        1 = -\frac{1}{2}(0) + C \implies C = 1 \\
        x_{s2}(t) = -\frac{1}{2}t + 1
    \end{gather*}
    Solving for the collision time $t_c$:
    \begin{gather*}
        -\frac{1}{2}t + 1 = \frac{1}{2}t - 1 \\
        -t + 2 = t - 2 \\
        2t = 4 \\
        t_c = 2 \\
        x_c = x_{s2}(2) = -\frac{1}{2}(2) + 1 = 0
    \end{gather*}
    For the third shock, taking the derivative of $x_{s3}$:
    \begin{gather*}
        u^- = 1,\ u^+ = -1 \\
        \frac{dx_{s3}}{dt} = \frac{\frac{1}{2}(-1)^2 - \frac{1}{2}(1)^2}{-1 - 1} = \frac{0}{-2} = 0 \\
        dx_{s3} = 0 \,dt \\
        \int \,dx_{s3} = \int 0 \,dt \\
        x_{s3}(t) = C
    \end{gather*}
    Using the point $(x_c, t_c) = (0, 2)$:
    \begin{gather*}
        C = x_c = 0 \\
        x_{s3}(t) = 0
    \end{gather*}
    Final piecewise solution:
    \begin{gather*}
        \text{For } t \leq 2 \\
        u(x,t) = \begin{cases}
                     1, & x < \frac{1}{2}t - 1 \\
                     0, & \frac{1}{2}t - 1 \leq x \leq -\frac{1}{2}t + 1 \\
                     -1, & x > -\frac{1}{2}t + 1
        \end{cases} \\\\
        \text{For } t > 2 \\
        u(x,t) = \begin{cases}
                     1, & x < 0 \\
                     -1, & x > 0
        \end{cases}
    \end{gather*}

    \newpage
    \question{
    \textbf{Feedback}\\
    How long did you spend on this homework? Was it easy/medium/hard? Do you have any feedback for the instructor/TA? Do you like the inclusion of bonus problems?
    }
    This homework took me a little while but that was because I was busy with work from other classes.
    I liked how the problems on this homework were mostly the same style.
    I have no opinion on the bonus questions since I normally don't have time to do them.
\end{document}